\chapter*{ABSTRAK}
\addcontentsline{toc}{chapter}{ABSTRAK}

{\fontsize{11}{13.2}\selectfont
Model proses yang dihasilkan algoritma \textit{process discovery} konvensional pada data pemeliharaan kerap rumit dan sukar dibaca, sebab deskripsi aktivitas ditulis bebas oleh teknisi sehingga memunculkan ratusan label berbeda untuk tindakan yang sebenarnya serupa. Penelitian ini memanfaatkan Generative AI, khususnya \textit{Large Language Model} (LLM), sebagai anotator untuk menstandarkan label aktivitas pada setiap \textit{event} melalui abstraksi semantik berbasis taksonomi dua belas kelas. Data berupa log pemeliharaan satu unit Pembangkit Listrik Tenaga Diesel (PLTD) diolah menjadi dua skenario, yaitu \textit{event log as-is} yang memakai \textit{raw labels} dan \textit{event log enriched} yang memakai label hasil anotasi LLM, lalu masing-masing dibangun modelnya menggunakan \textit{Inductive Miner} dan dievaluasi dengan empat metrik \textit{conformance}. Anotasi LLM menyeragamkan 112 aktivitas unik menjadi 12 kategori, memangkas jumlah \textit{place} sebesar 67,07\% dan \textit{transition} sebesar 70,09\%. Pengujian \textit{conformance} menunjukkan \textit{generalization} meningkat tajam dari 0,0670 menjadi 0,7267, sementara \textit{fitness} tetap tinggi pada 0,9590. Penurunan \textit{precision} dari 0,8821 menjadi 0,4729 mencerminkan berkurangnya \textit{overfitting} model \textit{as-is}, bukan penurunan kualitas. Hasil ini memperlihatkan bahwa abstraksi semantik berbasis LLM mampu mengubah model proses yang besar dan kusut menjadi model yang ringkas, dapat digeneralisasi, dan lebih bermakna secara operasional.

\vspace{1cm}

\noindent\textbf{Kata Kunci:} \textit{process mining}, \textit{process discovery}, \textit{event abstraction}, \textit{large language model}, abstraksi semantik, \textit{conformance checking}, pemeliharaan pembangkit listrik
}
